\textbf{- 参考资料：}
\href{https://e.math.cornell.edu/people/belk/numbertheory/CyclotomicPolynomials.pdf}{域与分圆多项式}~\cite{cyclotomic-polynomial}

\subsection{定义}
\label{subsec:cyclotomic-def}
\begin{tcolorbox}[title={\textbf{\tboxdef{\ref*{subsec:cyclotomic-def}} 分圆多项式}}]
 \textbf{第 $\bm{n}$ 个分圆多项式：}以所有 $n$ 次本原单位根为根的多项式，即：
 
 \[ \Phi_n(x) = \prod_{\zeta \in P(n)} (x - \zeta)  = \prod_{\substack{0 \leq k \leq n-1,\\ 
 \text{gcd}(k, n) = 1}} (x - \omega^k) \textcolor{white}{......} \text{，其中 } \omega = e^{2\pi i/n} \]

 回顾欧拉公式：$e^{2k\pi i/n} = \cos\left(\dfrac{2k\pi}{n}\right) + i\cdot\sin\left(\dfrac{2k\pi}{n}\right)$

$ $

一些预先计算出的分圆多项式如下：
 
\begin{multicols}{2}
$\Phi_1(x) = x - 1$
\newline $\Phi_2(x) = x + 1$
\newline $\Phi_3(x) = x^2 + x + 1$
\newline $\Phi_4(x) = x^2 + 1$
\newline $\Phi_5(x) = x^4 + x^3 + x^2 + x + 1$
\newline $\Phi_6(x) = x^2 - x + 1$
\newline $\Phi_7(x) = x^6 + x^5 + x^4 + x^3 + x^2 + x + 1$
\newline $\Phi_8(x) = x^4 + 1$
\newline $\Phi_9(x) = x^6 + x^3 + 1$
\newline $\Phi_{10}(x) = x^4 - x^3 + x^2 - x + 1$
\end{multicols}
\end{tcolorbox}


例如，

\[\Phi_4(x) = \prod_{\substack{0 \leq k \leq 3,\\ 
 \text{gcd}(k, 4) = 1}} (x - \omega^k) = (x - \omega^1)(x - \omega^3) = (x - e^{2\pi i / 4})(x - e^{2\cdot3\pi i / 4}) = (x - e^{\pi i / 2})(x - e^{3\pi i / 2})\]

$= \left(x - \left(\cos\left(\dfrac{\pi}{2}\right) + i \cdot \sin\left(\dfrac{\pi}{2}\right) \right) \right) \cdot \left(x - \left(\cos\left(\dfrac{3\pi}{2}\right) + i \cdot \sin\left(\dfrac{3\pi}{2}\right) \right) \right)$

$= (x - i)(x + i) = x^2 + 1$



\subsection{定理}
\label{subsec:cyclotomic-theorem}


\begin{tcolorbox}[title={\textbf{\tboxtheorem{\ref*{subsec:cyclotomic-theorem}.1} 第 $M$ 个分圆多项式的根}}]

假设 $M$ 是 2 的幂，并考虑第 $M$ 个分圆多项式 $\Phi_M(x) = x^n + 1$，其中 $M = 2n$。则第 $M$ 个分圆多项式的根为 $\omega, \omega^3, \omega^5, \cdots, \omega^{2n-1}$，其中 $\omega = e^{i\pi/n}$。

\end{tcolorbox}

\begin{proof}

根据 \autoref{subsec:cyclotomic-def} 中的定义~\ref*{subsec:cyclotomic-def}，$\Phi_M(x)$ 的根为 $e^{2k\pi i/M} = e^{2k\pi i/(2n)} = e^{k\pi i/n}$，其中 $0 \leq k < M = 2n$ 且 $\textsf{gcd}(k, M = 2n) = 1$，因此 $k = \{1, 3, 5, \cdots, 2n-1\}$。若令 $\omega = e^{i\pi/n}$，则 $\Phi_M(x)$ 的根为 $\omega, \omega^3, \omega^5, \cdots, \omega^{2n-1}$。

\end{proof}

\begin{tcolorbox}[title={\textbf{\tboxtheorem{\ref*{subsec:cyclotomic-theorem}.2} 将多项式分解为分圆多项式的乘积}}]
对于任意正整数 $n$，有 \[ x^n - 1 = \prod_{d \mid n} \Phi_d(x) \]
\end{tcolorbox}
\begin{myproof}
\begin{enumerate}
    \item $x^n - 1$ 的根是全部 $n$ 次单位根。因此，$x^n - 1 = (x - \omega^0)(x - \omega^1)...(x - \omega^{n-1})$，其中 $\zeta = \omega^k$。
    \item 定理~\ref*{subsec:roots-theorem}.3 表明，每个 $n$ 次单位根（$\omega^k$）都是某个 $d$ 次本原单位根，其中 $d$ 整除 $n$。换言之，每个 $n$ 次单位根都属于某个满足 $d \mid n$ 的 $P(d)$。另一方面，由定义有 $\Phi_d(x) = \Pi_{\zeta \in P(d)} (x - \zeta)$。因此，$x^n - 1$ 是所有满足 $d \mid n$ 的 $\Phi_d(x)$ 的乘积。
\end{enumerate}
\end{myproof}
\begin{tcolorbox}[title={\textbf{\tboxtheorem{\ref*{subsec:cyclotomic-theorem}.3} 分圆多项式的整数系数性质}}]
分圆多项式的所有系数都是整数。
\end{tcolorbox}
\begin{myproof}
对 $n$ 使用强归纳法。当 $n=1$ 时，$\Phi_1(x)=x-1$ 是首一整系数多项式。设 $n>1$，并假设结论对每个正整数 $d<n$ 成立，令
\[
F(x)=\prod_{\substack{d\mid n\\d<n}}\Phi_d(x).
\]
则 $F(x)$ 是 $\mathbb{Z}[x]$ 中的首一多项式。由定理~\ref*{subsec:cyclotomic-theorem}.2，在 $\mathbb{C}[x]$ 中有 $x^n-1=F(x)\Phi_n(x)$。以首一多项式 $F$ 作长除法只涉及整数运算，因此商和余式均属于 $\mathbb{Z}[x]$。由 $\mathbb{C}$ 上多项式除法的唯一性，商必为 $\Phi_n$，余式为零。因此 $\Phi_n\in\mathbb{Z}[x]$，归纳完成。
\end{myproof}
\begin{tcolorbox}[title={\textbf{\tboxtheorem{\ref*{subsec:cyclotomic-theorem}.4} $\bm{\Phi_{nk}(x)}$ 的公式}}]
若 $k \mid n$，则 $\Phi_{nk}(x) = \Phi_n(x^k)$。
\end{tcolorbox}
\begin{myproof}
对于满足 $k\mid n$ 的正整数，由定理~\ref*{subsec:order-theorem}.3，$\operatorname{ord}(\zeta)=nk$ 当且仅当 $\operatorname{ord}(\zeta^k)=n$。因此，$\Phi_n(x^k)$ 与 $\Phi_{nk}(x)$ 具有完全相同的根。这些根均为单根：$\Phi_n$ 的每个根 $\eta$ 都非零且为单根，而多项式 $x^k-\eta$ 的导数 $kx^{k-1}$ 在根处非零，所以它有 $k$ 个互异复根。两个多项式都首一且具有相同的单根，因此相等。映射 $\zeta\mapsto\zeta^k$ 在这些根集合上是 $k$ 对一，而非一对一；若对每个原像都乘入一次因子，就会错误地重复计算 $\Phi_n$ 的根。
\end{myproof}